\chapter{System Architecture}
\label{ch:architecture}

Chapter~\ref{ch:ensemble} presents prediction blending: two independent model branches combined at the forecast level.
This chapter presents two alternative ensemble architectures, feature stacking and residual stacking, and compares all three.


%% ──────────────────────────────────────────────
\section{Three Architectures}
\label{sec:three-architectures}
%% ──────────────────────────────────────────────

Figure~\ref{fig:architecture-comparison} shows the data and prediction flow for each architecture side by side.

\begin{figure}[htbp]
\centering
\resizebox{\textwidth}{!}{%
\begin{tikzpicture}[
  node distance=0.6cm and 0.3cm,
  >=Stealth,
  block/.style={draw, rounded corners, minimum height=0.7cm, minimum width=2.2cm,
    align=center, font=\scriptsize},
  data/.style={block, fill=blue!8},
  comp/.style={block, fill=green!8},
  model/.style={block, fill=orange!8},
  output/.style={block, fill=orange!12, font=\scriptsize\bfseries},
  arrow/.style={-{Stealth[length=2mm]}, thick},
  panellabel/.style={font=\small\bfseries, text=defblue},
]

% ── Panel A: Feature Stacking ──
\begin{scope}[xshift=0cm]
  \node[panellabel] at (1.5,5.2) {(A) Feature Stacking};
  \node[data] (a-seq) at (0,4) {E-mini L2\\Sequences};
  \node[model] (a-lstm) at (0,3) {LSTM};
  \node[comp] (a-emb) at (0,2) {Embedding\\($k$-dim)};
  \node[data] (a-tab) at (3,4) {Tabular\\Features (L0--7)};
  \node[comp] (a-concat) at (1.5,1) {Concatenate};
  \node[model] (a-lgbm) at (1.5,0) {LightGBM};
  \node[output] (a-out) at (1.5,-1) {Forecast};

  \draw[arrow] (a-seq) -- (a-lstm);
  \draw[arrow] (a-lstm) -- (a-emb);
  \draw[arrow] (a-emb) -- (a-concat);
  \draw[arrow] (a-tab) |- (a-concat);
  \draw[arrow] (a-concat) -- (a-lgbm);
  \draw[arrow] (a-lgbm) -- (a-out);
\end{scope}

% ── Panel B: Residual Stacking ──
\begin{scope}[xshift=7cm]
  \node[panellabel] at (1.5,5.2) {(B) Residual Stacking};
  \node[data] (b-tab) at (1.5,4) {Tabular\\Features};
  \node[model] (b-har) at (1.5,3) {HAR};
  \node[comp] (b-res1) at (1.5,2) {Residuals$_1$};
  \node[model] (b-lgbm) at (1.5,1) {LightGBM};
  \node[comp] (b-res2) at (1.5,0) {Residuals$_2$};
  \node[model, dashed] (b-lstm) at (1.5,-1) {LSTM (opt.)};
  \node[output] (b-out) at (1.5,-2) {$\sum$ Forecasts};

  \draw[arrow] (b-tab) -- (b-har);
  \draw[arrow] (b-har) -- (b-res1);
  \draw[arrow] (b-res1) -- (b-lgbm);
  \draw[arrow] (b-lgbm) -- (b-res2);
  \draw[arrow, dashed] (b-res2) -- (b-lstm);
  \draw[arrow] (b-lstm) -- (b-out);
  \draw[arrow, gray] (b-tab.east) -- ++(0.8,0) |- (b-lgbm.east);
  % E-mini sequences feed optional LSTM
  \node[data] (b-seq) at (-1.8,-1) {E-mini\\Sequences};
  \draw[arrow, dashed] (b-seq) -- (b-lstm);
\end{scope}

% ── Panel C: Prediction Blending (Ch.11) ──
\begin{scope}[xshift=13cm]
  \node[panellabel] at (1.5,5.2) {(C) Prediction Blending};
  \node[data] (c-tab) at (0,4) {Tabular\\Features};
  \node[data] (c-seq) at (3,4) {E-mini\\Sequences};
  \node[model] (c-lgbm) at (0,2.5) {LightGBM};
  \node[model] (c-lstm) at (3,2.5) {LSTM};
  \node[comp] (c-blend) at (1.5,1) {Weighted\\Average};
  \node[output] (c-out) at (1.5,0) {Forecast};

  \draw[arrow] (c-tab) -- (c-lgbm);
  \draw[arrow] (c-seq) -- (c-lstm);
  \draw[arrow] (c-lgbm) -- (c-blend);
  \draw[arrow] (c-lstm) -- (c-blend);
  \draw[arrow] (c-blend) -- (c-out);
  \node[font=\tiny, text=gray, anchor=north] at (1.5,-0.5) {(Ch.~\ref{ch:ensemble})};
\end{scope}

\end{tikzpicture}%
}
\caption{Three ensemble architectures.
(A)~LSTM embedding concatenated with tabular features for a single LightGBM.
(B)~Each model trains on residuals from the prior stage; final forecast is the sum.
(C)~Independent models blended at prediction level (Chapter~\ref{ch:ensemble}).}
\label{fig:architecture-comparison}
\end{figure}


%% ──────────────────────────────────────────────
\section{Feature Stacking}
\label{sec:feature-stacking}
%% ──────────────────────────────────────────────

The LSTM processes E-mini L2 5-min bars and LOB features (78 time steps per day) and produces a $k$-dimensional embedding vector (default $k=32$, or $k=1$ for a scalar forecast).
This embedding is concatenated with the ${\sim}$80--120 tabular features from Layers~0--7 to form a single expanded feature set.
LightGBM then trains on the combined input.

The approach has one fundamental problem: LightGBM cannot back-propagate gradients through the LSTM, so the embedding is never optimized for the tree's $\QLIKE$ objective.
The LSTM learns representations that minimize its own loss, which may not be what the tree needs.

\begin{table}[htbp]
\centering
\caption{Feature stacking: pros and cons.}
\label{tab:feature-stacking}
\small
\begin{tabular}{@{}lp{10cm}@{}}
\toprule
\textbf{Pros} & Single training pass; LSTM learns representations the tree can exploit \\
\midrule
\textbf{Cons} & Gradient isolation (tree cannot backprop into LSTM); embedding not optimized for tree objective; debugging harder; no RV literature demonstrates this beating alternatives \\
\bottomrule
\end{tabular}
\end{table}


%% ──────────────────────────────────────────────
\section{Residual Stacking}
\label{sec:residual-stacking}
%% ──────────────────────────────────────────────

Stage~1: HAR baseline (OLS) produces a forecast and residuals.
Stage~2: LightGBM trains on Stage~1 residuals with the full tabular feature set.
Stage~3 (optional): LSTM trains on Stage~2 residuals from E-mini sequences.
The final forecast is the sum of all stage forecasts.

Each model specializes by construction.
HAR captures the multi-scale autoregressive structure that dominates at all horizons.
LightGBM captures nonlinear patterns the HAR misses (regime interactions, jump-asymmetry effects).
The LSTM, if used, captures whatever sequential dynamics remain in the residuals.

\begin{table}[htbp]
\centering
\caption{Residual stacking: pros and cons.}
\label{tab:residual-stacking}
\small
\begin{tabular}{@{}lp{10cm}@{}}
\toprule
\textbf{Pros} & Each model has a distinct role; no gradient isolation; clean residual targets; aligns with HARQ-X direction; supported by recent RV literature \\
\midrule
\textbf{Cons} & Sequential training (each stage depends on prior); residual signal may be weak at later stages \\
\bottomrule
\end{tabular}
\end{table}


%% ──────────────────────────────────────────────
\section{Comparison}
\label{sec:architecture-comparison}
%% ──────────────────────────────────────────────

\begin{table}[htbp]
\centering
\caption{Three-way architecture comparison.}
\label{tab:architecture-comparison}
\small
\begin{tabular}{@{}lp{3.2cm}p{3.2cm}p{3.8cm}@{}}
\toprule
\textbf{Dimension} & \textbf{Feature Stacking} & \textbf{Residual Stacking} & \textbf{Pred.\ Blending (Ch.~\ref{ch:ensemble})} \\
\midrule
Complexity &
  High (joint training) &
  Moderate (sequential) &
  Low (independent) \\[6pt]
Gradient flow &
  Broken (tree cannot backprop into LSTM) &
  Clean (residual targets) &
  N/A (independent) \\[6pt]
Literature &
  Weak (no RV paper) &
  Strong (HARQ-X, recent lit.) &
  Strong \citep{Optiver2021, Bucci2020} \\[6pt]
Fallback &
  Must retrain tree without embedding &
  Drop Stage~3; keep HAR + LightGBM &
  Drop one model; keep the other \\[6pt]
Interpretability &
  Opaque (embedding) &
  Clear (stage contributions) &
  Clear (individual forecasts) \\
\bottomrule
\end{tabular}
\end{table}


%% ──────────────────────────────────────────────
%% Boxes
%% ──────────────────────────────────────────────

\begin{keyidea}[Residual Stacking Gives Each Model a Distinct Role]
HAR captures multi-scale $\RV$ persistence.
LightGBM captures nonlinear patterns the HAR misses.
LSTM (if used) captures whatever regime dynamics remain.
Each model trains on residuals from the prior stage, so roles are distinct by construction.
\end{keyidea}

\begin{warning}[No Evidence for Feature Stacking in RV Forecasting]
No paper in the RV literature demonstrates LSTM-embedding-to-GBDT feature stacking beating prediction blending or residual stacking at any forecast horizon.
The gradient isolation problem (Chapter~\ref{ch:ensemble}) compounds this: embeddings are never optimized for the tree objective.
\end{warning}
